% From mitthesis package
% Documentation: https://ctan.org/pkg/mitthesis

\ProvidesFile{mydesign.tex}[2026/01/31 v1.07 Typographic design options for thesis]

% ==> you can change this file however you want. The content below is just for illustration. <==


%%%%%%%%%  Change page margins  %%%%%%%%%%%%%%%%%%%%%%%%%%%%
% 
% The default thesis margin is 1 inch all around. You may want different margins (e.g., to add a gutter for binding),
% in which case you can use the \newgeometry command from the geometry package.  Refer to the package documentation
% for details.
%
% mitthesis defaults: [top=1in,bottom=1in,left=1in,right=1in,marginparwidth=50pt,headsep=12pt,footskip=0.5in]
%
% The following tells the geometry package to use a two-sided layout with a 1 cm binding offset on the inside 
% 	and 1 inch margins all around, reducing textwidth slightly (by 0.7 cm). See geometry documentation, Section 8.2.
%
%\newgeometry{twoside, bindingoffset=1cm,margin=1in,marginparwidth=50pt,headsep=12pt,footskip=0.5in}


%%%%%%%%%%  Color support  %%%%%%%%%%%%%%%%%%%%%%%%%%%%%%%%%

%% Color package: xcolor. 
%% Change this if you prefer something else

\usepackage[dvipsnames,svgnames,x11names]{xcolor}

%% xcolor defines a broad set of color names that can be invoked in your document (see xcolor documentation).
%% can add option [table] to xcolor to use color in tables 

%%%%%%%%%%  Hyperlink and line number colors  %%%%%%%%%%%%%%

\AtBeginDocument{
    % Using color names from xcolor package
    \hypersetup{
    	linkcolor=Blue3,
    	citecolor=Blue3,
    	urlcolor=violet,
    	filecolor=red, 
    %	anchorcolor=yellow,%  not all pdf viewers recognize this field (although Firefox does): hyperref issues a warning, which can be ignored
    %	colorscheme=phelype,% overrides link, cite, url, file colors with a preset scheme, through \DocumentMetadata
    	}
    %
    \ifmit@lineno
	    \renewcommand{\linenumberfont}{\sffamily\mdseries\tiny\color{violet}}% line numbers will be sans-serif, medium weight, tiny, and violet
	\fi
}

%%%%%%%%%  Caption support  %%%%%%%%%%%%%%%%%%%%%%%%%%%%%%%%
% 

\IfFormatAtLeastTF{2025/11/01}
{
    %% If your latex format is at least 2025/11/01, you should use this approach:
    \NewSocketPlug{caption/label}{plug-A}{{\bfseries #1:\space}} % This strictly addresses the caption label, not caption text.
    \AssignSocketPlug{caption/label}{plug-A}
    %
    \AddToHook{cmd/@makecaption/before}{\small} % This styles the entire caption, but label socket above will override it.
    %
    % These commands make caption label bold with a colon separating label from text;
    %    \small affects both label and text (unless you add a sizing command to plug-A).
}{
    %% For older formats, use this. See documentation of caption package for details.
    \RequirePackage{caption}
    %% This sets caption labels in bold face type with a colon separator. 
    	\captionsetup[figure]{labelfont={bf},labelsep=colon} % change to boldface label with colon separator
    	\captionsetup[table]{labelfont={bf}, labelsep=colon} % change to boldface label with colon separator
}

%%%%%%%%%  Customize titles and section headings  %%%%%%%%%%

% To directly modify section headings as follows, see https://latexref.xyz/dev/latex2e.html#g_t_005c_0040startsection
%
% The following commands adds \raggedright to section headings (compare to report.cls)
%
% NB: for fonts with a predefined bold \mathversion, the  >> fontset <<  file inserts \mathversion{bold} into these commands.
%
\renewcommand\section{\@startsection {section}{1}{\z@}%
                                   {-3.5ex \@plus -1ex \@minus -.2ex}%
                                   {2.3ex \@plus.2ex}%
                                   {\mathversion{bold}\normalfont\Large\bfseries\raggedright}}
\renewcommand\subsection{\@startsection{subsection}{2}{\z@}%
                                   {-3.25ex\@plus -1ex \@minus -.2ex}%
                                   {1.5ex \@plus .2ex}%
                                   {\mathversion{bold}\normalfont\large\bfseries\raggedright}}
\renewcommand\subsubsection{\@startsection{subsubsection}{3}{\z@}%
                                   {-3.25ex\@plus -1ex \@minus -.2ex}%
                                   {1.5ex \@plus .2ex}%
                                   {\normalfont\normalsize\bfseries}}
\renewcommand\paragraph{\@startsection{paragraph}{4}{\z@}%
                                   {3.25ex \@plus1ex \@minus.2ex}%
                                   {-1em}%
                                   {\normalfont\normalsize\bfseries}}
\renewcommand\subparagraph{\@startsection{subparagraph}{5}{\parindent}%
                                   {3.25ex \@plus1ex \@minus .2ex}%
                                   {-1em}%
                                   {\normalfont\normalsize\bfseries}}

%% This makes numbered chapter headings san-serif
%\patchcmd{\@makechapterhead}{\Huge}{\Huge\sffamily}{}{}
%\patchcmd{\@makechapterhead}{\huge}{\huge\sffamily}{}{}

%% This makes UNnumbered chapter headings san-serif
%\patchcmd{\@makeschapterhead}{\Huge}{\Huge\sffamily}{}{}

%%%%%%%%%%%%

%% The titlesec package provides broader customizations of titles and headings - refer to its documentation
%% NB: titlesec may break pdf tagging !!
%
%\RequirePackage{titlesec}
%
%% these titlesec commands just automate the bold math in some headings (rest follows default styles)
%
%\titleformat{\section}{\mathversion{bold}\normalfont\Large\bfseries}{\thesection }{1em}{}
%\titleformat{\subsection}{\mathversion{bold}\normalfont\large\bfseries}{\thesubsection}{1em}{}
%\titleformat{\subsubsection}{\mathversion{bold}\normalfont\normalsize\bfseries}{\thesubsubsection}{1em}{}


%%%%%%%%%  Customize list environments  %%%%%%%%%%%%%%%%%%%%
% 
% See documentation of enumitem package for details of how to customize lists
% NB: enumitem may break pdf tagging !!

%\RequirePackage{enumitem}
